\section{Bayesian Data Quality Model}

\label{sec:quality}
% ══════════════════════════════════════════════════════════════════════════════

\subsection{Observation Reliability}

Not all observations are equally reliable. We model the reliability
$q_{obs}$ of an observation $o$ as:

\begin{equation}
  q_{obs} = p(\text{correct ID} \mid \text{observer}, \text{AI},
  \text{evidence}, \text{context})
\end{equation}

\subsection{Generative Model}

We define a Bayesian generative model:

\begin{enumerate}
  \item True species identity $z \sim p(z \mid \text{location}, \text{date})$
    drawn from the local species pool.
  \item Observer identification $y_h \sim p(y_h \mid z, \text{skill}_h)$
    dependent on true identity and observer skill.
  \item AI identification $y_{ai} \sim p(y_{ai} \mid z, \text{image})$
    dependent on true identity and image quality.
  \item Community votes $\{v_1, \ldots, v_C\}$ where
    $v_c \sim p(v_c \mid z, \text{skill}_c)$.
\end{enumerate}

The posterior over true identity is:
\begin{equation}
  p(z \mid y_h, y_{ai}, \{v_c\}) \propto
  p(z) \cdot p(y_h \mid z) \cdot p(y_{ai} \mid z) \cdot
  \prod_c p(v_c \mid z)
\end{equation}

\subsection{Observer Skill Estimation}

Observer skill $\text{skill}_h$ is modeled as a latent variable updated
through the observation history. We use a beta-binomial model:
\begin{equation}
  \text{skill}_h \sim \text{Beta}(\alpha_h, \beta_h)
\end{equation}
where $\alpha_h$ and $\beta_h$ are updated with each confirmed
correct or incorrect identification, respectively. Species-group-specific
skill parameters allow for differential expertise (e.g., expert birder,
novice botanist).

\subsection{Quality Scores in Practice}

\begin{table}[H]
\centering
\caption{Distribution of observation quality scores across the
ZooPlatform dataset.}
\label{tab:quality}
\begin{tabular}{lcc}
\toprule
\textbf{Quality tier} & \textbf{$q_{obs}$ range} & \textbf{Percentage} \\
\midrule
Research grade & $\geq 0.95$ & 61.3\% \\
Needs review & $0.80$--$0.95$ & 23.7\% \\
Casual & $0.60$--$0.80$ & 11.4\% \\
Low confidence & $< 0.60$ & 3.6\% \\
\bottomrule
\end{tabular}
\end{table}

61.3\% of observations achieve research-grade quality scores ($q_{obs}
\geq 0.95$), comparable to professional survey data. The quality model
enables downstream analyses to appropriately weight observations rather
than applying binary include/exclude filters.


% ══════════════════════════════════════════════════════════════════════════════
